\section*{Hoofdstuk \ref{ch:g11:seq} --- Rijen: een eerste kennismaking}

\begin{solution}{exo:g11:seq:1}
$u_n = \frac{n}{n+1}$: $u_1 = \frac12$, $u_2 = \frac23$, $u_3 = \frac34$.

$u_0 = 5$, $u_{n+1} = 3u_n - 2$: $u_1 = 13$, $u_2 = 37$, $u_3 = 109$.

$u_n = (-1)^n n$: $u_1 = -1$, $u_2 = 2$, $u_3 = -3$.
\end{solution}

\begin{solution}{exo:g11:seq:2}
$u_{10} = 7 + 10 \times (-3) = -23$ en $u_{25} = 7 - 75 = -68$.

Voor $(v_n)$: $v_8 = v_3 + 5d$ geeft $26 = 11 + 5d$, dus $d = 3$; daarna is
$v_0 = v_3 - 3d = 11 - 9 = 2$.
\end{solution}

\begin{solution}{exo:g11:seq:3}
$u_8 = 5 \times 2^8 = 1280$.

Voor $(v_n)$: $v_4 = v_2\, q^2$ geeft $48 = 12 q^2$, dus $q^2 = 4$ en
$q = 2$ (de termen zijn positief). Daarna is
$v_0 = \frac{v_2}{q^2} = \frac{12}{4} = 3$.
\end{solution}

\begin{solution}{exo:g11:seq:4}
$1 + \dots + 500 = \frac{500 \times 501}{2} = 125\,250$.

$4 + 7 + \dots + 61$ is \omterm{def:g11:seq:arithmetic}{rekenkundig} met $d = 3$ en
$\frac{61 - 4}{3} + 1 = 20$ termen: som $20 \times \frac{4 + 61}{2} = 650$.

$1 + \frac12 + \dots + \frac{1}{2^{10}}$ is \omterm{def:g11:seq:geometric}{meetkundig} met $q = \frac12$ en
$11$ termen:
$\frac{1 - (1/2)^{11}}{1 - 1/2} = 2\left(1 - \frac{1}{2048}\right)
= \frac{2047}{1024}$.
\end{solution}

\begin{solution}{exo:g11:seq:5}
$u_{n+1} - u_n = 4(n+1) - 1 - 4n + 1 = 4$: \omterm{def:g11:seq:arithmetic}{rekenkundig} met $d = 4$.

$\frac{v_{n+1}}{v_n} = \frac{2^{n+1}}{3^{n+2}} \cdot \frac{3^{n+1}}{2^n}
= \frac23$: \omterm{def:g11:seq:geometric}{meetkundig} met $q = \frac23$.

$w_0 = 0$, $w_1 = 2$, $w_2 = 6$: de verschillen $2$ en $4$ zijn ongelijk, dus
niet \omterm{def:g11:seq:arithmetic}{rekenkundig}; $\frac{w_1}{w_0}$ is niet eens gedefinieerd, en de
verhoudingen $\frac{w_2}{w_1} = 3 \neq \frac{w_3}{w_2} = 2$: geen van beide.
\end{solution}

\begin{solution}{exo:g11:seq:6}
De aantallen stoelen per rij vormen een \omterm{def:g11:seq:arithmetic}{rekenkundige rij}: eerste term $16$,
\omterm{def:g11:seq:arithmetic}{verschil} $2$. De laatste (twintigste) rij telt $16 + 19 \times 2 = 54$
stoelen. Het totaal is $20 \times \frac{16 + 54}{2} = 700$ stoelen.
\end{solution}

\begin{solution}{exo:g11:seq:7}
Na $n$ uur telt de populatie $500 \times 2^n$. Om 20 uur is $n = 8$:
$500 \times 256 = 128\,000$ bacteriën. We hebben
$500 \times 2^n > 10^6$ nodig, dus $2^n > 2000$: omdat $2^{10} = 1024$ en
$2^{11} = 2048$, overtreft de populatie één miljoen voor het eerst na $11$
volle uren, om 23 uur.
\end{solution}

\begin{solution}{exo:g11:seq:8}
$c_2 = 1.002 \times 100 + 100 = 200.20$ en
$c_3 = 1.002 \times 200.20 + 100 \approx 300.60$. De verschillen
$c_2 - c_1 = 100.20$ en $c_3 - c_2 \approx 100.40$ zijn niet gelijk, dus is
$(c_n)$ niet \omterm{def:g11:seq:arithmetic}{rekenkundig}; de verhoudingen $\frac{c_2}{c_1} = 2.002$ en
$\frac{c_3}{c_2} \approx 1.50$ zijn evenmin gelijk, dus is ze ook niet
\omterm{def:g11:seq:geometric}{meetkundig}. (Gemengde recursies van het type ``vermenigvuldig en tel dan
op'' los je op met de hulprijtruc van \cref{exo:g11:seq:11}.)
\end{solution}

\begin{solution}{exo:g11:seq:9}
$u_{n+1} - u_n = (n+1)^2 - 8(n+1) - n^2 + 8n = 2n - 7$: negatief voor
$n \leq 3$, positief voor $n \geq 4$. Dus daalt $(u_n)$ tot
$u_4 = 16 - 32 = -16$ en stijgt ze daarna: ze is niet monotoon.

$(v_n)$ heeft positieve termen en
\[
\frac{v_{n+1}}{v_n} = \frac{3^{n+1}}{(n+1)!} \cdot \frac{n!}{3^n}
= \frac{3}{n+1},
\]
wat $> 1$ is voor $n \leq 1$, $= 1$ voor $n = 2$ en $< 1$ voor $n \geq 3$:
de \omterm{def:g11:seq:sequence}{rij} stijgt tot $v_2 = v_3 = \frac92$ en daalt daarna.
\end{solution}

\begin{solution}{exo:g11:seq:10}
De eerste $n$ termen zijn $u_0, \dots, u_{n-1}$, met $u_0 = 3$ en
$u_{n-1} = 3 + 4(n-1) = 4n - 1$. Hun som is
\[
n \times \frac{3 + (4n-1)}{2} = n(2n + 1) = 903,
\]
dus $2n^2 + n - 903 = 0$. Hier is $\Delta = 1 + 4 \times 2 \times 903 =
7225 = 85^2$, en $n = \frac{-1 + 85}{4} = 21$ (de negatieve \omterm{def:g11:quad:discriminant}{wortel} valt
weg). Controle: $21 \times 43 = 903$.
\end{solution}

\begin{solution}{exo:g11:seq:11}
\emph{1.} $u_1 = 5$, $u_2 = 9$, $u_3 = 17$: elke term is één meer dan $4$,
$8$, $16$, wat $u_n = 2^{n+1} + 1$ doet vermoeden.

\emph{2.} Met $v_n = u_n - 1$:
\[
v_{n+1} = u_{n+1} - 1 = 2u_n - 1 - 1 = 2(u_n - 1) = 2v_n,
\]
dus is $(v_n)$ \omterm{def:g11:seq:geometric}{meetkundig} met \omterm{def:g11:seq:geometric}{reden} $2$ en eerste term $v_0 = u_0 - 1 = 2$.

\emph{3.} Bijgevolg is $v_n = 2 \times 2^n = 2^{n+1}$ en
$u_n = v_n + 1 = 2^{n+1} + 1$, wat het vermoeden bevestigt. (Het getal $1$
dat in $v_n$ wordt afgetrokken is het vaste punt van $x \mapsto 2x - 1$;
hetzelfde idee duikt voor $u_{n+1} = au_n + b$ opnieuw op in
\cref{ch:g12:seq}.)
\end{solution}

\begin{solution}{pb:g11:seq:1}
\textbf{1.} $h_1 = 1$, $h_2 = 3$, $h_3 = 7$.

\textbf{2.} Om de grootste schijf te verplaatsen moeten de $n$ schijven
erboven eerst naar de reservepen verhuizen ($h_n$ zetten); de grote schijf
steekt over ($1$ zet); en de $n$ schijven moeten er daarna weer bovenop
klimmen ($h_n$ zetten): $h_{n+1} = 2h_n + 1$.

\textbf{3.} $v_{n+1} = h_{n+1} + 1 = 2h_n + 2 = 2v_n$: \omterm{def:g11:seq:geometric}{meetkundig} met \omterm{def:g11:seq:geometric}{reden}
$2$ en $v_1 = 2$, dus $v_n = 2^n$ en $h_n = 2^n - 1$.

\textbf{4.} $2^{64} - 1 \approx 1.8 \times 10^{19}$ seconden; gedeeld door
$3 \times 10^7$ seconden per jaar geeft dat ongeveer $6 \times 10^{11}$
jaar --- zeshonderd miljard jaar, veertig keer de ouderdom van het heelal.
De monniken mogen gerust koffiepauzes nemen.

\textbf{5.} Bekijk in elke geldige oplossing de eerste zet van de onderste
schijf: op dat ogenblik moeten de andere $n$ schijven alle op de ene
overblijvende pen liggen (minstens $h_n$ zetten om ze daar te krijgen), en na
de laatste zet van de onderste schijf moeten ze er alle weer bovenop komen
(minstens $h_n$ zetten meer): elke oplossing heeft dus minstens $2h_n + 1$
zetten nodig. De recursie is een ondergrens én een bovengrens: $2^n - 1$ is
optimaal.

\textbf{6.} $1, 1, 2, 3, 5, 8, 13, 21, 34, 55, 89, 144$.

\textbf{7.} Niet \omterm{def:g11:seq:arithmetic}{rekenkundig} ($2 - 1 = 1$ en $3 - 2 = 1$, maar
$5 - 3 = 2$: de verschillen veranderen); niet \omterm{def:g11:seq:geometric}{meetkundig} ($\frac21 = 2$
maar $\frac32 = 1.5$). \omterm{def:g11:func:monotone}{Stijgend}: voor $n \geq 2$ is
$F_{n+1} - F_n = F_{n-1} > 0$.

\textbf{8.} $F_k = F_{k+2} - F_{k+1}$, dus
\[
\sum_{k=1}^{n} F_k = (F_3 - F_2) + (F_4 - F_3) + \dots +
(F_{n+2} - F_{n+1}) = F_{n+2} - F_2 = F_{n+2} - 1 .
\]
Voor $n = 6$: $1 + 1 + 2 + 3 + 5 + 8 = 20 = F_8 - 1 = 21 - 1$.

\textbf{9.} $F_k F_{k+1} - F_{k-1} F_k = F_k (F_{k+1} - F_{k-1}) =
F_k \cdot F_k = F_k^2$; optellen telescopeert tot $F_n F_{n+1} - F_1 F_0$
(met $F_0 = 0$): de som van de kwadraten is $F_n F_{n+1}$. Voor $n = 4$:
$1 + 1 + 4 + 9 = 15 = F_4 F_5 = 3 \times 5$.

\textbf{10.} $F_5 F_3 - F_4^2 = 5 \times 2 - 9 = 1$;
$F_6 F_4 - F_5^2 = 8 \times 3 - 25 = -1$;
$F_7 F_5 - F_6^2 = 13 \times 5 - 64 = 1$: afwisselend $\pm 1$. Dat \omterm{def:g11:seq:arithmetic}{verschil}
van één tussen $F_{n+1} F_{n-1}$ en $F_n^2$ is precies de gewonnen of
verloren vierkante eenheid van de goochelaar: een vierkant
$F_n \times F_n$ in stukken snijden en tot een rechthoek
$F_{n+1} \times F_{n-1}$ hersamenstellen moet één eenheid scheppen of
opslokken --- de spleet.

\textbf{11.} $F_{n+2} = F_{n+1} + F_n \geq F_n + F_n = 2F_n$ (de \omterm{def:g11:seq:sequence}{rij}
stijgt): om de twee indices minstens een verdubbeling --- groei die minstens
\omterm{def:g11:seq:geometric}{meetkundig} is met \omterm{def:g11:seq:geometric}{reden} $\sqrt2$ per index.

\textbf{12.} $1.5$; $1.667$; $1.6$; $1.625$; $1.615$; $1.619$; $1.618$;
$1.618$. Als $r_n \to L$: deel $F_{n+2} = F_{n+1} + F_n$ door $F_{n+1}$, wat
$r_{n+1} = 1 + \frac{1}{r_n}$ geeft, dus $L = 1 + \frac1L$, dus
$L^2 = L + 1$: $L = \varphi = \frac{1 + \sqrt5}{2}$, de gulden snede van
\cref{pb:g10:algebra:1}. De konijnen vermenigvuldigen zich in goud.

\textbf{13.} $v_{n+1} = u_{n+1} - \ell = a u_n + b - \ell$; omdat
$\ell = a\ell + b$, is dat $a(u_n - \ell) = a v_n$: \omterm{def:g11:seq:geometric}{meetkundig} met \omterm{def:g11:seq:geometric}{reden}
$a$. Bijgevolg is $v_n = a^n v_0$ en $u_n = a^n (u_0 - \ell) + \ell$.

\textbf{14.} \omterm{pb:g10:functions:1}{Vast punt}: $\ell = 1.01\ell - 300$ geeft $\ell = 30\,000$. Dus
$d_n = 1.01^n (10\,000 - 30\,000) + 30\,000 = 30\,000 - 20\,000 \times
1.01^n$.

\textbf{15.} $d_n \leq 0$ vereist $1.01^n \geq 1.5$:
$1.01^{40} \approx 1.489$ en $1.01^{41} \approx 1.504$: de $41$ste
aflossing wist de schuld uit (en is iets kleiner dan $300$). Totaal
terugbetaald: net onder $41 \times 300 = 12\,300$ euro --- de geleende
$10\,000$ kostte ongeveer $2\,300$ euro rente.

\textbf{16.} \omterm{pb:g10:functions:1}{Vast punt} $\ell = \frac{1000}{1 - 1.02} = -50\,000$, dus
$p_n = 1.02^n \times 100\,000 - 50\,000$. Na $10$ jaar is
$1.02^{10} \approx 1.219$: $p_{10} \approx 71\,900$ inwoners.

\textbf{17.} $\frac{1000 \times 1001}{2} = 500\,500$; en
$2^{20} - 1 = 1\,048\,575$.

\textbf{18.} Van $7$ tot $502$ met stappen van $5$:
$\frac{502 - 7}{5} + 1 = 100$ termen; som
$= 100 \times \frac{7 + 502}{2} = 25\,450$.

\textbf{19.} Saldo $= 100 \times \frac{1.005^{60} - 1}{1.005 - 1} \approx
100 \times \frac{0.3489}{0.005} \approx 6\,977$ euro --- waarvan $6\,000$
gestort en ongeveer $977$ verdiend: meetkundige sommen zijn de moedertaal
van de bank.

\textbf{20.} Expliciete formules beantwoorden ``wat is $u_{1000}$'' meteen;
recursies beschrijven hoe systemen werkelijk evolueren --- de kunst bestaat
erin het tweede in het eerste om te zetten. \omterm{def:g11:seq:arithmetic}{Rekenkundige rijen} tellen op,
meetkundige vermenigvuldigen, en elke familie heeft haar eigen somformule
(de koppeling van Gauss; de verdubbelingstruc). De truc met het vaste punt
en de hulprij zet elke affiene recursie om in een meetkundige --- leningen,
bevolkingen en de toren bezweken er alle voor. Fibonacci gehoorzaamt aan geen
van beide families, en toch vingen telescoperende identiteiten haar sommen en
kwadraten; haar volledige portret (een exacte formule, de gulden limiet)
wacht op sterker gereedschap.
\end{solution}
